\section{The characteristic-two family}\label{sec:characteristic-two}

We turn to the ordered factorization \((u,v)=(2,1)\).  Thus
\[
 q=8m,\qquad n=6m+1,\qquad \lambda=2.
\]
Near equality write
\begin{equation}\label{eq:characteristic-two-asymptotic}
 k=32m^2+cm+o(m),\qquad T=32m+h+o(1).
\end{equation}
The integer pair condition and the external incidence count give
\begin{equation}\label{eq:characteristic-two-linear}
 6c-h\ge76,\qquad
 e=2(c+h-8)m+o(m),
\end{equation}
where \(e=\sum_{Q\notin\mathcal K}(\sigma_n(Q)-2)\) is the total excess
above the required two \(n\)-secants.

The dual set \(\mathcal L^*\) has only \(O(q)\) odd-secants.  For sufficiently
large even \(q\), this is below the Sz\H{o}nyi--Weiner threshold
\((\lfloor\sqrt q\rfloor+1)(q+1-\lfloor\sqrt q\rfloor)\).  Their stability
theorem~\cite[Theorem~1.1]{SWeven} therefore writes it as the symmetric difference
\begin{equation}\label{eq:even-type-repair}
 \mathcal L^*=\mathcal E\mathbin{\triangle}R,
\end{equation}
where every line meets \(\mathcal E\) in an even number of points and
\(|R|=r=\lceil\delta/(q+1)\rceil=O(1)\), with \(\delta\) the number of
odd-secants of \(\mathcal L^*\).

\begin{lemma}[The even-type case]\label{lem:even-type-case}
\coverage{absent}
Let \(\mathcal K\) be a \((k,n)\)-arc with parameters
\eqref{eq:characteristic-two-asymptotic}, suppose every external point lies on
at least two \(n\)-secants, and let \(\mathcal L^*\) be the dual point set of
all its \(n\)-secants.  If \(R=\varnothing\) in
\eqref{eq:even-type-repair}, then \(c\ge16\).
\end{lemma}

\begin{proof}
The balanced external pair is \(\{2,3\}\).  Put \(D=6k-nT\), the signed
internal deficiency from six.  The incidence equations and the unrestricted
integer pair bound give
\begin{align*}
 D&=(6c-6h-32)m+o(m),\\
 e&=2(c+h-8)m+o(m),\\
 S_{\mathrm{aff}}&=(12c-2h-152)m+o(m).
\end{align*}
If \(D\ge0\), the balanced internal pair is \(\{5,6\}\), and
\(S_{\mathrm{aff}}\) is the available unrestricted pair slack to first
order.  If \(D<0\), the balanced pair is \(\{6,7\}\); convexity of the
internal \(\Phimin\) term makes the same low-branch affine continuation
\(S_{\mathrm{aff}}\) an upper bound on the actual unrestricted slack.
If every multiplicity is even, restricting the internal balanced distribution
to even integers raises its minimum by \(|D|/2\), while the external
restriction raises its minimum by \(e/2\).  The total increase is
\((|D|+e)/2\).  Since \(|D|\ge D\), it is at least
\[
 \frac{D+e}{2}=(4c-2h-24)m+o(m).
\]
In either branch, feasibility forces \(S_{\mathrm{aff}}\) to be at least this
large, which simplifies to \(c\ge16\).
\end{proof}

\begin{lemma}[Parity of the repair]\label{lem:repair-parity}
\coverage{absent}
Suppose \(R\ne\varnothing\) in \eqref{eq:even-type-repair}.  If \(r\) is
odd, then
\begin{equation}\label{eq:odd-repair-bound}
 c+h\ge9.
\end{equation}
If \(r\) is even, then
\begin{equation}\label{eq:even-repair-bound}
 c+h\ge10.
\end{equation}
Moreover, \(r\equiv T\equiv h\pmod2\).
\end{lemma}

\begin{proof}
Every even-type set has even cardinality, so
\(r\equiv|\mathcal L^*|=T\pmod2\); the leading term \(32m\) is even, giving
the last congruence.

Dualize the points of \(R\) to \(r\) lines of the original plane.  On each
such line, at most \(r-1\) points can cancel with the other generators in the
symmetric difference.  Thus at least \(q-r+2\) points have odd
\(n\)-secant multiplicity.  At most \(n=3q/4+1\) of them lie in
\(\mathcal K\), so a single generator contributes at least
\(q/4-r+1\) external points to the excess \(e\).

If \(r\) is odd, then \(r\ge1\) and hence \(e\ge q/4-O(1)=2m-O(1)\).
If \(r\) is even, then \(r\ge2\); taking the union of two generator lines
loses only their intersection, and \(e\ge q/2-O(1)=4m-O(1)\).  Comparing
with \eqref{eq:characteristic-two-linear} proves
\eqref{eq:odd-repair-bound} and \eqref{eq:even-repair-bound}.
\end{proof}

\begin{proof}[Proof of Theorem~\ref{thm:characteristic-two}]
Suppose the theorem fails.  Pass to a sequence for which the coefficient
\(c_m=(k-32m^2)/m\) is bounded above by
\(73/6-8\varepsilon\).  The classical incidence bound bounds \(c_m\) below,
so pass to a subsequence with \(c_m\to c\).  The leading incidence bound and
Lemma~\ref{lem:secant-number-near-equality} give
\(T-32m=O(1)\).  Pass to a further subsequence with fixed integer difference
\(h\).

Equation~\eqref{eq:characteristic-two-linear} and \(c\le73/6\) give
\[
 6c-6h-32\ge424-30c\ge59.
\]
Thus \(D>0\) for all sufficiently large \(m\), so the balanced internal pair
in this contradiction range is \(\{5,6\}\).

As in the characteristic-three proof, the exact decomposition of \(-F(T)\)
into the two nonnegative balancing losses and
Equation~\eqref{eq:balancing-loss} show that only \(O(m)\) internal degrees
lie outside \(\{5,6\}\), and only \(O(m)\) external degrees lie outside
\(\{2,3\}\).  Moreover,
\[
 D=6k-nT=\sum_{P\in\mathcal K}(6-\rho_n(P))=O(m).
\]
For \(z\notin\{5,6\}\), the absolute value of \(6-z\) is at most twice the
corresponding summand in \eqref{eq:balancing-loss}; hence the number of
internal degree-five points is \(O(m)\).  Since every external degree is at
least two and
\(e=\sum_Q(\sigma_n(Q)-2)=O(m)\), the number of external degree-three points
is also \(O(m)\).  These are the wrong-parity adjacent degrees, so only
\(O(m)=O(q)\) dual lines have the wrong parity.  Therefore
Sz\H{o}nyi--Weiner's theorem gives \eqref{eq:even-type-repair} with bounded
integer \(r\).  Pass to a further subsequence on which \(r\) is fixed.
Lemma~\ref{lem:even-type-case} excludes \(r=0\) in the range under
consideration.  If \(h\) is odd, Lemma~\ref{lem:repair-parity} and
\eqref{eq:characteristic-two-linear} give
\[
 c\ge\max\left\{\frac{76+h}{6},9-h\right\}.
\]
The minimum over odd \(h\) is \(73/6\), attained at \(h=-3\).  If \(h\) is
even, the corresponding bound is
\[
 c\ge\max\left\{\frac{76+h}{6},10-h\right\},
\]
whose minimum over even \(h\) is \(37/3\), attained at \(h=-2\).  Both
contradict the chosen sequence.  Since \(q=8m\), this proves
\eqref{eq:characteristic-two-bound}.
\imports{SWeven:repair}
\uses{lem:even-type-case,lem:repair-parity,lem:secant-number-near-equality}
\end{proof}

The parity argument uses only the repair theorem and the fact that an
\(n\)-secant contains at most \(n\) points of the arc.  It does not require a
classification of even-type sets of size \(4q+O(1)\).
